\title{Byte Latent Transformer: Patches Scale Better Than Tokens}

\begin{abstract}
We present the Byte Latent Transformer ({\textsc BLT}), a novel large language model (LLM) design operating directly at the byte level, which is the first of its kind to achieve performance parity with tokenization-based LLMs at large scales, while delivering notable gains in both inference speed and resilience.
{\textsc BLT} processes input by grouping bytes into adaptively sized patches, treating these as the fundamental units for computation. The segmentation of patches relies on estimating the entropy of upcoming bytes, dedicating increased model attention and compute resources to more complex or information-rich data regions. We conduct the inaugural scaling analysis of byte-level models with flop-constrained budgets, spanning model sizes up to 8B parameters and encompassing 4T training bytes. Our experiments establish that training models on raw bytes—eliminating any reliance on a predetermined vocabulary—is a viable approach for scaling. Enhanced efficiencies in both training and inference are achieved through the dynamic use of extended patches when sequences are more predictable, complemented by improvements in both reasoning tasks and generalization to rare scenarios. Summarizing, at a fixed inference computation budget, {\textsc BLT} achieves superior scaling compared to token-based architectures by increasing both patch sizes and model capacities simultaneously.
\end{abstract}

\section{Introduction}
\label{section:intro}

We present the Byte Latent Transformer~(\textbf{BLT}), an architecture that dispenses with tokenization by directly processing raw byte sequences. For the first time, this approach achieves parity with tokenization-dependent models at scale, while delivering notable gains in both computational efficiency and model robustness (\S\ref{sec:robustness}).

Although most large language models (LLMs) operate in a nearly end-to-end manner, they universally rely on tokenization as a heuristic preprocessing step to partition bytes into a predetermined vocabulary of tokens. This process introduces biases into string compression and presents a series of limitations, including sensitivity to the particular domain or modality~\citep{dagangetting}, increased susceptibility to perturbations in input~(\S\ref{sec:robustness}), insufficient awareness of orthographic variation~\citep{edman2024cute}, and inequitable performance across languages~\citep{liang2023xlm,petrov2024language,limisiewicz2024myte}.

Tokenization has traditionally been a crucial step because training large language models (llms) directly on byte-level inputs remains prohibitively expensive at scale, mainly due to the substantial increase in sequence length~\citep{xue2022byt5}. To address this, earlier studies have introduced more computationally efficient self-attention mechanisms~\citep{el2020characterbert, clark2022canine} or even adopted architectures without attention~\citep{wang2024mambabyte}~(\S\ref{section:related_work}). Nevertheless, these strategies are chiefly effective for training \textit{small models}. When scaling up, the bulk of a Transformer’s computation is consumed by the extensive feed-forward network layers applied to every byte, whereas the attention mechanism contributes relatively little to the total cost.

In order to optimize computational resource distribution, we introduce an adaptive and trainable technique that organizes bytes into \textit{patches}~(\S\ref{section:patching}), accompanied by a novel architecture capable of integrating both byte-level and patch-level data. Distinct from tokenization approaches, {\textsc BLT} operates without a predetermined patch vocabulary. Instead, flexible sequences of bytes are projected into latent patch spaces using compact, trainable encoder and decoder components. Our empirical findings demonstrate that this framework enables a \textit{greater} computational efficiency compared to traditional tokenization-based systems.

Traditional tokenization-based language models assign equal computational resources to each token. This approach prioritizes uniformity but sacrifices efficiency, as token boundaries are established by compression heuristics that do not necessarily correspond to the difficulty of making predictions. The foundation of our proposed architecture lies in the principle that computation should be adaptively allocated according to task demands. For instance, generating the final characters of most words generally requires less modeling capacity, since these are straightforward, low-entropy predictions compared to selecting an initial word at the beginning of a sentence. This motivation is evident in {\textsc BLT}'s design~(\S\ref{section:architecture}), which employs three transformer modules: a pair of compact byte-level \textit{local models} and a larger, more powerful global \textit{latent transformer}~(\autoref{fig:arch_high_level}). To facilitate adaptive computation, {\textsc BLT} organizes bytes into variable-length patches by analyzing the entropy associated with predicting the next byte. This results in context-dependent groupings of bytes, each with similar information content.

This work introduces the first comprehensive scaling investigation of byte-level models regulated by floating point operation counts, extending up to 8B parameters and spanning 4T training bytes. Our results demonstrate that it is feasible to train large-scale, end-to-end models directly from raw bytes, eliminating the need for a predetermined vocabulary during tokenization. In terms of training efficiency measured by flop usage,\footnote{We quantify a model's computational requirements via the total count of Floating Point OPerations (flops).} {\textsc BLT} achieves parity with Llama 3, but it realizes as much as a 50\% reduction in inference-time flops compared to conventional models (\S\ref{section:scaling}). 

Furthermore, we find that modeling at the byte level yields substantial gains when dealing with infrequent or idiosyncratic data patterns. The {\textsc BLT} architecture demonstrates increased robustness to corrupted input and exhibits superior performance in character-level linguistic tasks—including tests of orthography, phonology, and scenarios with limited parallel data for machine translation (\S\ref{sec:robustness}).

Moreover, our approach enables concurrent expansion of both model and patch sizes without exceeding a fixed inference-time computational budget. Utilizing longer patch sizes generally reduces computation, freeing resources for further growth of the global latent transformer, since it is activated less frequently. Through flop-constrained scaling experiments at inference (\autoref{fig:fixed_inference_scaling}), we observe that {\textsc BLT} displays markedly better scaling properties compared to standard, tokenization-dependent architectures.

In conclusion, the primary contributions of this work are as follows: 1) We present {\textsc BLT}, a byte-level latent language model architecture that adaptively allocates computation resources, thereby enhancing flop efficiency, 2) We demonstrate that our approach attains training flop-matched performance comparable to Llama 3 up to the 8B parameter scale, with the flexibility to exchange modest reductions in evaluation scores for up to a 50\% increase in flop efficiency, 3) The {\textsc BLT} architecture introduces a novel approach to scaling llms, enabling expansion in model size while adhering to a fixed inference computational cost, 4) Our results highlight the superior robustness of {\textsc BLT} models to noisy inputs and their sensitivity to sub-word information in input sequences, which traditional token-based llms often overlook. The training and inference code for {\textsc BLT} is publicly available at~\url{https://github.com/facebookresearch/blt}.

\section{Patching: From Individual Bytes to Groups of Bytes}
\label{section:patching}

Dividing a byte sequence into distinct \textit{patches} enables {\textsc BLT} to adaptively assign computational resources depending on contextual requirements. As illustrated in Figure~\ref{fig:patching_types}, there exist multiple strategies for forming patches from a stream of bytes. More specifically, a patching function $f_p$ partitions a byte sequence $\pmb{x} = \{x_i, |i=1, \ldots, n\}$ of length $n$ into $m < n$ discrete patches $\pmb{p} = \{p_j | j = 1, \ldots, m\}$ by mapping each element $x_i$ to an indicator in \{0,1\}, where a value of 1 signifies the initiation of a new patch. For both token-based systems and those organized around patches, the primary computational burden is shaped by how many steps the central Transformer processes. For {\textsc BLT}, this correlates to the count of patches the patching function dictates are necessary for representing the data. As a result, the predominant factor affecting processing cost, whether during model training or inference for any particular patching function~(\S\ref{section:flops}), is the mean \textit{patch size}. In the following sections, we describe three methods for patching: using a fixed quantity of bytes per patch~(\S\ref{section:static-patch}), employing whitespace as patch boundaries~(\S\ref{section:white-patch}), and adopting a dynamic approach where entropy estimates from a small byte language model guide patch segmentation~(\S\ref{section:dyn-patch}). We also examine incremental patching strategies and elaborate on distinctions between patching and tokenization~(\S\ref{section:bpe-patch}).

\subsection{Strided Patching Every K Bytes}
\label{section:static-patch}

A common and simple approach to dividing bytes is by segmenting them into patches of a constant size $k$, as seen in MegaByte~\citep{yu2023megabyte}. This method, relying on a uniform stride, allows for ease of implementation during both training and inference, provides an uncomplicated means to adjust the average size of a patch, and grants direct control over computational resources expended (flops). Nonetheless, this approach has notable drawbacks. Chief among them is the inflexibility in computation allocation: computational effort may be squandered on transformer step $j$ when only low-importance elements, such as whitespace in source code, are processed; conversely, inadequate computation may be allotted to highly informative segments, such as mathematical expressions. In addition, this fixed patching scheme leads to fragmented and context-insensitive tokenization of similar byte strings—for instance, occurrences of the same word may be split inconsistently.

\subsection{Space Patching}
\label{section:white-patch}

\citet{slagle2024spacebyte} introduces a straightforward yet impactful refinement to strided patching, where new patches are initiated following any space-like bytes\footnote{
    Space-like bytes refer to any byte that is neither a Latin character, a digit, nor a UTF-8 continuation byte. Additionally, each patch must contain at least one byte that is not space-like. }—positions that commonly correspond to the boundaries of linguistic units in numerous languages. With Space patching, each word receives a dedicated step within the latent transformer (entailing increased computational cost), ensuring a consistent patching scheme for words across different sequences, and allocating computation specifically to challenging prediction points, which often coincide with spaces. For instance, generating the initial byte in a response to the prompt ``Who composed the Magic Flute?\uline{\hspace{.6em}}'' tends to be substantially more complex than predicting subsequent bytes such as those in ``Mozart,'' since the set of plausible candidates is drastically reduced after the first character, making the remainder easier to infer. Nevertheless, space patching faces limitations: it may not robustly accommodate all languages and domains, and is fundamentally constrained by its inability to adjust patch size dynamically. To address these challenges, we next propose a novel patching strategy built around the observation that the prediction of initial bytes in words is characteristically difficult, while also furnishing an intrinsic way to control patch sizing.

\subsection{Entropy Patching: Using Next-Byte Entropies from a Small Byte LM}
\label{section:dyn-patch}

In place of traditional rule-based methods like using whitespace, we adopt a data-driven methodology for detecting positions of elevated uncertainty in next-byte predictions. We propose \textit{entropy patching}, a technique that determines patch boundaries based on entropy calculations.

To this end, a compact byte-level auto-regressive language model is trained on the {\textsc BLT} training data. Next-byte entropies are then computed according to the language model distribution $p_{e}$ across the byte vocabulary $\mathcal{V}$:

\begin{equation}
H(x_i) = \sum_{v \in \mathcal{V}} p_{e}(x_i=v|\pmb{x}_{<i}) \log p_{e}(x_i=v|\pmb{x}_{<i})
\end{equation}

We employ two approaches for detecting patch boundaries based on the entropies $H(x_i)$. The initial method selects points that exceed a fixed, global entropy threshold, demonstrated in~\autoref{fig:patching}. The alternative technique focuses on locating points where the entropy significantly increases compared to the preceding value. This latter method can be viewed as highlighting locations where the pattern of nearly monotonic entropy decrease within the patch is disrupted.

\begin{align*}
    \text{Global Constraint}\hspace{1cm}H(x_t)& > \theta_g\\
    \text{Approx. Monotonic Constraint}\hspace{1cm}H(x_t)& - H(x_{t-1}) > \theta_r
\end{align*}

Patch boundaries are determined through a simple preprocessing procedure performed at the dataloading stage. This approach contrasts with that of~\citet{nawrot-etal-2023-efficient}, in which a classifier is trained to infer entropy-based patch boundaries. In our experiments~(\S\ref{section:experiments}), we evaluate both strategies for discriminating between bytes of low and high entropy.

\subsection{The Byte-Pair Encoding (BPE) Tokenizer and Incremental Patching}
\label{section:incremental-patching}
\label{section:bpe-patch}

A large number of state-of-the-art llms, such as the baseline Llama 3 examined in this work, rely on subword tokenization methods like bpe~\citep{gage1994new, sennrich-etal-2016-neural}. In this context, we define ``tokens'' as sequences of bytes selected from a \textit{finite} set of possible options established before training begins. This stands in contrast to ``patches,'' which are formed dynamically and do not rely on a predetermined vocabulary. Notably, an essential distinction is that token-based models lack direct interaction with the raw byte-level features, unlike when using patches.

An important advancement offered by {\textsc BLT} in comparison to tokenization-based approaches is its reconfiguration of the relationship between vocabulary size and computational requirements. Conventional llms experience a trade off wherein expanding the vocabulary leads to longer average token lengths, which in turn reduces the total number of model steps but simultaneously enlarges the final projection layer's output dimension. This inherent constraint limits the capacity of tokenization-based models to meaningfully adjust token length and inference overhead. As an illustrative case, Llama 3 extends the average token length from 3.7 to 4.4 bytes, yet does so by quadrupling the size of its embedding matrix relative to Llama 2.

During generation, {\textsc BLT} must determine at each step of the byte sequence whether it is currently at a patch boundary, as this dictates whether to engage additional computation through the Latent Transformer. Importantly, this decision is made without knowledge of subsequent parts of the sequence that have not yet been produced. As a consequence, patching strategies cannot rely on access to future bytes when segmenting the sequence. More formally, a patching scheme $f_p$ is said to be incrementally patchable if it adheres to the following criterion:
$$f_p(\pmb{x}_{<i}) = f_p(\pmb{x})_{<i}$$
The bpe algorithm does not meet the requirements of an incremental patching scheme, since the same prefix may be segmented differently depending on how the rest of the sequence unfolds. Thus, it fails to satisfy the above condition\footnote{Incorporating a unique delimiter to mark patch boundaries can convert bpe into an incremental scheme, but at the cost of increasing sequence length.}.

\section{{\textsc BLT} Architecture}
\label{section:architecture}
The {\textsc BLT} framework consists of a substantial global autoregressive language model that processes patch-level representations, complemented by two more lightweight, local models responsible for mapping sequences of bytes into patches and reconstructing byte sequences from patch representations, as illustrated in~\autoref{fig:arch_high_level}.

\subsection{Latent Global Transformer Model}

\textit{The Latent Global Transformer} refers to an autoregressive transformer model $\mathcal{G}$ composed of $l_{\mathcal{G}}$ layers, which translates a series of latent input patch embeddings $p_j$ into corresponding output patch embeddings $o_j$. Here, the notation $j$ is consistently reserved for indexing patches, while $i$ is used for bytes. The global architecture employs a block-causal attention mask~\citep{dubey2024llama}, limiting each token's attention to patches up to and including its own within a document. As the primary consumer of floating-point operations during both pre-training and inference, careful scheduling of this model allows adaptive allocation of computational resources based on the complexity encountered in various segments of the input and output.

\subsection{Local Encoder}
\label{section:local}

\textit{The Local Encoder Model}, represented as $\mathcal{E}$, is a compact transformer variant comprised of $l_{\mathcal{E}} << l_{\mathcal{G}}$ layers, designed primarily for efficient transformation of an input byte sequence $b_i$ into rich patch-level representations $p_j$. Unlike standard transformer architectures, this model integrates a cross-attention layer after each transformer block, facilitating the aggregation of byte embeddings into patch representations (see~\autoref{fig:crossattn}).  
Initially, each byte $b_i$ is converted into an embedding $x_i$ via a matrix in $\mathbb{R}^{256 \times h_{\mathcal{E}}}$. Optional enhancements to these embeddings can be incorporated through hash-embeddings~(\S\ref{sec:hash-emb}).  
These representations undergo transformation through multiple pairs of transformer and cross-attention layers~(\S\ref{sec:enc-cross-attn}), yielding patch embeddings $p_i$ for subsequent processing by the global transformer $\mathcal{G}$.  
Within the transformer layers, a \textit{local block causal} attention pattern is applied; specifically, each byte attends to a window of $w_{\mathcal{E}}$ preceding bytes. While this attention window may span dynamic patch boundaries, it does not extend across document boundaries. The next sections provide further explanation regarding the embeddings and the structure of the cross-attention component.

\subsubsection{Encoder Hash n-gram Embeddings}
\label{sec:ngram-emb}
\label{sec:hash-emb}
Capturing strong and informative representations at every timestep $i$ critically depends on embedding knowledge of the preceding sequence of bytes. Within {\textsc BLT}, this is accomplished by representing the current byte $b_i$ not only in isolation but also in conjunction with its surrounding bytes as a byte n-gram. Accordingly, at each position $i$, we begin by forming byte-grams

\begin{align}
    g_{i,n}=\{b_{i-n + 1},\ldots, b_i\}
\end{align}

for each byte position $i$ and $n$ from three to eight.\footnote{
We omit byte-grams of size $n$ or more when $i<n$. }

Subsequently, we present hash $n$-gram embeddings, where every byte $n$-gram is projected by a hash function onto a fixed-size index within an embedding table $E_{n}^{hash}$, for each $n \in\{3,4,5,6,7,8\}$~\citep{bai2010learning}. The embedding produced in this way is summed with the original byte embedding before being normalized, after which it serves as input to the local encoder. The composite embedding is then computed as follows:

\begin{align}
e_i &= x_i + \sum_{n = 3,...,8}E_{n}^{hash}(\text{Hash}(g_{i,n})) \\
\text{where, Hash}(g_{i,n}) &= \text{RollPolyHash}(g_{i,n}) \% |E_{n}^{hash}|
\end{align}

We adjust $e_i$ by dividing it by the total number of $n$-gram sizes plus one and implement the RollPolyHash approach as specified in Appendix~\ref{appendix:rollpolyhash}. Section~\ref{section:ablations} presents an ablation study analyzing how varying the parameters $n$ and embedding table size for $n$-gram hash embeddings influences trends observed in flop-controlled scaling laws. Alongside the use of hash $n$-gram embeddings, we also explored frequency-based $n$-gram embeddings; further information regarding this line of investigation is given in Appendix~\ref{appendix:ngrams}.

\subsubsection{Encoder Multi-Headed Cross-Attention}
\label{sec:enc-cross-attn}
Our approach draws heavily from the input cross-attention mechanism used in the Perceiver architecture~\citep{jaegle2021perceiver}, with the primary modification that here, the latent vectors are associated with specific patch representations, rather than being a pre-defined, fixed set~(\autoref{fig:crossattn}). Additionally, attention is limited to those bytes constituting each individual patch. In this scheme, a query vector is created for every patch $p_j$ by applying a pooling operation over the byte-level embeddings within $p_j$, and subsequently projecting via a linear transformation, $\mathcal{E}_{C} \in \mathbb{R}^{h_{\mathcal{E}} \times (h_{\mathcal{E}} \times U_{\mathcal{E}})}$, where $U_{\mathcal{E}}$ specifies the encoder’s cross-attention head count. Explicitly, with $f_{\text{bytes}}(p_j)$ representing the sequence of bytes for the patch $p_j$, the following computations are performed:

\begin{align}
P_{0,j} &= \mathcal{E}_{C}(f_\text{bytes}((p_j)), f~ \text{denotes the pooling operation}) \\
P_l &= P_{l-1} + W_o\left(\text{softmax}\left(\frac{QK^T}{\sqrt{d_k}}\right)V\right) \\
\text{with } Q_j &= W_q(P_{l-1,j}),~ K_i = W_k(h_{l-1,i}),~ V_i = W_v(h_{l-1,i}) \\
h_l &= \text{Encoder-Transformer-Layer}_l(h_{l-1})
\end{align}

Here, $P \in \mathbb{R}^{n_p \times h_{\mathcal{G}}}$ denotes the representations for $n_p$ patches, which serve as the inputs to the global model. The initialization of these patch representations is achieved by pooling the byte embeddings $e_i$ that make up each patch $p_j$. The matrices $W_q$, $W_k$, $W_v$, and $W_o$ are linear transformations applied to generate queries, keys, values, and outputs, respectively, with keys and values derived from the byte-level hidden states $h_i$ of the preceding layer ($e_i$ for the initial layer). A patch-specific masking technique is employed such that the query for patch $Q_j$ interacts solely with the keys and values associated with the bytes within the same patch $j$. Since we employ multi-head attention for the $Q$, $K$, and $V$ representations, and because patch representations usually possess higher dimensionality ($h_{\mathcal{G}}$) compared to $h_{\mathcal{E}}$, we structure $P_l$ as a collection of attention heads, each of dimension $h_{\mathcal{E}}$ for the cross-attention step, before concatenating these heads to reconstruct the $h_{\mathcal{G}}$ dimension. A pre-LayerNorm is applied to queries, keys, and values, and no positional encodings are incorporated within this cross-attention mechanism. A residual connection is introduced around the cross-attention module to facilitate training.

\subsection{Local Decoder} 

In line with the architecture of the local encoder, the local decoder $\mathcal{D}$ utilizes a lightweight transformer structure comprising $l_{\mathcal{D}}$ layers, where $l_{\mathcal{D}} << l_{\mathcal{G}}$. This component translates a sequence of global patch embeddings, $o_j$, back into the original byte-level data, $y_i$. The decoder generates the byte sequence autoregressively, conditioned on previously generated bytes, drawing on the hidden states derived from the local encoder for the current sequence. It is organized into $l_{\mathcal{D}}$ repetitions of cross-attention followed by standard transformer layers. Notably, each decoding step begins with a cross-attention mechanism that establishes byte-level representations from the patch representations prior to the application of the transformer layer, which subsequently processes these reconstructed byte representations.

\subsubsection{Decoder Multi-headed Cross-Attention} 

Within the decoder's cross-attention mechanism, the typical configuration of queries and keys/values is inverted: here, the byte-representations act as queries, while the patch representations are designated as the keys and values. To initiate the cross-attention computation, the byte-representations are seeded with the byte embeddings derived from the final encoder layer, denoted as $h_{l_{\mathcal{E}}}$. For each subsequent layer $l$, the byte-representations $d_{l,i}$ are updated via the following procedure:

\begin{align}
D_0 &= h_{l_{\mathcal{E}}} \\
B_l &= D_{l-1} + W_o\left(\text{softmax}\left(\frac{QK^T}{\sqrt{d_k}}\right)V\right), \\
  &\text{where } Q_i = W_q(d_{l-1,i}),\quad K_i = W_k(\mathcal{D}_C(o_j)),\quad V_i = W_v(\mathcal{D}_C(o_j)) \\
  D_l &= \text{Decoder-Transformer-layer}_l(B_l)
\end{align}

In this formulation, $W_k$ and $W_v$ serve as key and value projection matrices that perform linear transformations on, and are partitioned by, the $\mathcal{D}_C$ splitting operation, which is applied to the global model's output patch representations $o_j$. The query projection matrix $W_q$ transforms the byte-level representations $d_{l-1}$ coming from the previous decoder transformer layer, or from $h_{l_{\mathcal{E}}}$ in the initial layer. The output projection matrix is denoted by $W_o$. Consequently, $B$ has shape $B \in \mathbb{R}^{h_{\mathcal{D}} \times n_b}$, with $n_b$ representing the total number of output bytes. The subsequent decoder layer outputs $D_l$ are obtained by passing the cross-attention block's output $B$ through a decoder transformer layer. Consistent with the approach taken in the local encoder cross-attention mechanism, our setup employs multiple attention heads, utilizes pre-layer normalization, omits positional embeddings, and includes a residual connection enclosing the cross-attention module.

\section{Experimental Setup}
\label{section:experiments}
We methodically construct tightly controlled experiments in order to evaluate {\textsc BLT} against models that use tokenization, ensuring that {\textsc BLT} is not unintentionally advantaged, such as by benefiting from longer sequence windows.

\subsection{Pre-training Datasets} In our experiments, all model variants undergo pre-training on two primary datasets: 1) The Llama 2 dataset~\citep{touvron2023llama}, which contains 2 trillion tokens sourced from a wide range of publicly accessible collections, systematically cleaned and filtered to enhance quality; and 2) {\textsc BLT}-1T: a newly constructed corpus consisting of 1 trillion tokens drawn from assorted public sources, also incorporating a segment of the pre-training material made available by Datacomp-LM~\citep{li2024datacomp}. The Llama 2 dataset underpins scaling law investigations to identify the optimal token count as suggested in~\cite{dubey2024llama}, thereby guiding architectural decisions for {\textsc BLT}. Meanwhile, {\textsc BLT}-1T supports comprehensive pre-training to facilitate direct comparisons with Llama 3 on downstream benchmark tasks. Notably, neither dataset contains information obtained from Meta’s products or services. In addition, baseline studies involving tokenizer-based models utilize the Llama 3 tokenizer—which possesses a vocabulary of 128K tokens—since it consistently outperformed the Llama 2 tokenizer across our evaluations.

\subsection{Entropy Model} For our experiments, the entropy model utilized is a byte-level language model that is trained using the same dataset as the full {\textsc BLT} system. By default—unless specified otherwise—we employ a transformer architecture consisting of 100M parameters, featuring 14 layers, each with a hidden dimensionality of 512, and leveraging sliding window attention over 512 bytes. All other hyperparameters align with those applied in our local and global transformer models. We explored several configurations, including variations in model size, receptive field length, and architecture as detailed in \autoref{section:ablations}. Notably, if the receptive field of this model is sufficiently limited, the resulting trained entropy model is amenable to representation as a compact lookup table.

\subsection{Entropy Threshold and Equalizing Context Length} For models that implement entropy-driven patching, we determine a patching threshold corresponding to a targeted average \textit{patch size} over the mixture of pretraining data. In {\textsc BLT}, in contrast to standard tokenization approaches, the \textit{patch size} is not constrained and can be selected freely, which has considerable effects on the context window accessible to the model. To ensure fairness by preserving the average context length and to prevent models with larger patch sizes from receiving an undue benefit, we regulate the number of bytes per batch so it stays approximately constant. As a result, models with larger patch sizes have their sequence lengths decreased. For experiments using Llama 2 data, we fix the context to 8k bytes, whereas for the {\textsc BLT}-1T dataset, we expand the context size to an average of 16k bytes, still adhering to an average batch size of 16M bytes.

Although the mean batch size remains unchanged, dynamic patching approaches result in variable byte-to-patch ratios during batch loading. To improve computational efficiency, our {\textsc BLT} training setup organizes batches by aggregating patches, thereby eliminating the need for additional padding operations in the costly latent transformer stage. This design guarantees a consistent patch count across all batches. For training purposes, we pad or, when necessary, truncate byte sequences to fixed lengths of 12k bytes for Llama 2 and 24k bytes for the {\textsc BLT}-1T datasets, which helps prevent excessive memory usage caused by sequences that include exceptionally large patches.

\subsection{Entropy Model Context}
\label{section:entropy-context}
Our empirical observations indicate that applying entropy patching tends to produce increasingly larger patches in tasks characterized by structured, repetitive content, such as multiple choice formats (refer to patching in an MMLU example in \autoref{fig:mmlu_patching}). This trend arises due to reduced entropy within the recurring elements present in the entropy model's context. Consequently, for the large-scale implementation of {\textsc BLT}-Entropy with a patch size of 4.5, we reinitialize the entropy context using line breaks and adopt the approximate monotonicity constraint, since it is less vulnerable to "entropy drift" associated with fluctuations in context length. Importantly, this modification influences only our entropy computations, and our method for determining the entropy threshold remains unchanged.

\subsection{FLOPs Estimation} 
\label{section:flops}

Our methodology for calculating transformer flops primarily draws from Chinchilla~\citep{hoffmann2022training}, which accounts for the flops involved in the feed-forward networks, the qkvo projections within self-attention modules, as well as the attention score and output projection computations. However, our approach diverges by treating the input embedding layer as an efficient lookup rather than a dense matrix operation, and thus we attribute zero flops to this step. Consistent with earlier research, we approximate that the backward pass incurs twice the computational cost, in terms of flops, compared to the forward pass.

In order to determine the number of flops \textit{per byte} for {\textsc BLT} models, we aggregate the computational costs associated with the local encoder transformer, the global latent transformer, the local decoder transformer, as well as the flops incurred by the cross attention components in both the encoder and decoder. This is expressed as:

\begin{align}
    \text{FL}_{\text{{\textsc BLT}}} &= \text{Transf. FL}(h_{\mathcal{G}}, l_{\mathcal{G}}, m=n_{ctx}/n_p, V=0) / n_p\\
    &+ \text{Transf. FL}(h_{\mathcal{E}}, l_{\mathcal{E}}, m=w_{\mathcal{E}}, V=0) \\
    &+ \text{Transf. FL}(h_{\mathcal{D}}, l_{\mathcal{D}}, m=w_{\mathcal{D}}, V=256) \\ 
    &+ \text{Cross Attn. FL}(h_{\mathcal{E}}, l_{\mathcal{E}}, m=n_p, r=n_p/k) \times k/n_p \\ 
    &+ \text{Cross Attn. FL}(h_{\mathcal{D}}, l_{\mathcal{D}}, m=k, r=k/n_p) 
\end{align}

Here, $n_{ctx}$ denotes the sequence length measured in bytes, while $n_p$ indicates the patch size. The parameter $r$ captures the proportion of query tokens relative to key/value tokens, and $k$ represents the ratio between the patch dimension and the byte dimension—that is, the count of local model splits that are concatenated to yield the global model representation ($k=2$ as illustrated in \autoref{fig:crossattn}). $V$ stands for the output projection's vocabulary size, utilized solely within the local decoder module. The applicable context length $m$ in the attention mechanism depends on whether the module operates over byte sequences or patch sequences. We accordingly adapt the attention flops calculations for every respective component. Precise formulations for both Transformer-FLOPs and Cross-Attention FLOPs computations can be found in Appendix~\ref{section:flops-appendix}.

\subsection{Bits-Per-Byte Estimation} Perplexity is meaningful only with a specific tokenizer, since it quantifies uncertainty over individual tokens. Therefore, to enable comparison between models operating at the byte level versus those at the token level, we adopt the Bits-Per-Byte (BPB) metric, as recommended in earlier research~\citep{xue2022byt5, yu2023megabyte, wang2024mambabyte}. BPB offers a measure of perplexity that does not depend on any particular tokenizer. More precisely:

\begin{align}
    \text{BPB}(x)&= \frac{\mathcal{L}_{CE}(\pmb{x})}{\ln(2)\cdot n_{\text{bytes}}}
\end{align}

where the uncertainty over the data $\pmb{x}$ as measured by the sum of the cross-entropy loss is normalized by the total number of bytes in $\pmb{x}$ and a constant. 

\subsection{Transformer Architecture Hyperparameters} In the design of all transformer blocks within {\textsc BLT}—covering both the local and global components—we base our architecture closely on Llama~3~\citep{dubey2024llama}. Specifically, we implement the SwiGLU activation function~\citep{swiglu} in the feed-forward networks, apply rotary positional embeddings (RoPE)~\citep{RoPe} configured with $\theta=500000$~\citep{xiong2024effective} exclusively in self-attention modules, and adopt RMSNorm~\citep{RMSNorm} for layer normalization. Self-attention layers using standard attention masks—such as \textit{block causal} or \textit{fixed-window block causal} masks—make use of Flash attention~\citep{dao2022flashattention}; in these layers, attention masks with a fixed window employ a width of 512. For the cross-attention layers, which rely on dynamic patch-dependent attention masks, we utilize Flex Attention\footnote{\url{https://pytorch.org/blog/flexattention}}. This allows us to generate efficient fused kernel implementations, considerably accelerating the training process.

\subsection{{\textsc BLT}-Specific Hyperparameters} In order to evaluate the performance of {\textsc BLT} models, we carry out experiments in two main areas: examining scaling behaviors and assessing downstream tasks. For these purposes, we examine models with varying parameter counts: 400M, 1B, 2B, 4B, and 8B. Detailed architectural hyperparameters for each model variant can be found in Appendix~Table~\ref{tab:arch}.

To initialize queries for the first cross-attention layer in the local encoder, we employ max-pooling. Across all {\textsc BLT} model sizes, we use $500,000$ hashes with a single hash function and set the n-gram length between 3 and 8. A fixed learning rate of $4e-4$ is adopted universally for these models. This learning rate aligns with the optimal value established via a hyperparameter sweep ranging from $1e-3$ to $1e-4$ at the 400M and 1B model sizes, revealing that equivalent settings are optimal for both token-based and {\textsc BLT} models. For scaling experiments using Llama-2 data, we adhere to the batch-size guidance from~\cite{dubey2024llama} or its byte-wise equivalent.

For optimization, the AdamW optimizer~\citep{adamw} is utilized, with $\beta_1=0.9$, $\beta_2=0.95$, and $\epsilon=10^{-8}$. We apply a linear warm-up over the first 2000 steps, followed by cosine decay of the learning rate down to zero. Additionally, we employ weight decay of 0.1 and enforce a global gradient clipping threshold of 1.0.

\section{Scaling Trends}
\label{section:scaling}

We provide a comprehensive overview of the scaling patterns observed in byte-level models, offering insights that can guide the continued scaling of {\textsc BLT} models. Our scaling analysis seeks to overcome the constraints found in earlier work on byte-level architectures through the following approaches: (a) Examining compute-optimal training conditions for comparative trends, (b) Training comparable 8B parameter models using substantial training data volumes (up to 1T tokens or 4T bytes) followed by downstream task evaluation, and (c) Assessing scaling behaviors under scenarios with fixed inference costs. In subsequent sections, we further explore particular benefits associated with processing byte sequences.

\subsection{Parameter Matched Compute Optimal Scaling Trends}
\label{section:parameter-matched-scaling}
Leveraging the Llama 2 dataset, we conduct training of multiple \textit{compute-optimal} bpe and {\textsc BLT} models at four different scales, with parameter counts ranging between 1B and 8B. Subsequently, we analyze how training flops correlate with language modeling performance using a representative fraction of the overall data mixture. The bpe models are trained according to the optimal parameter-to-data ratio specified by Llama~3~\citep{dubey2024llama}. This \textit{compute-optimal} regimen is theoretically constructed to maximize training dataset performance within a fixed computational budget~\citep{hoffmann2022training}, thus serving as a solid reference point. For every bpe model, an equivalent {\textsc BLT} model is also trained on the identical dataset, incorporating a Latent Transformer with the same parameter count and architecture as its bpe Transformer counterpart.

As shown in~\autoref{fig:scaling-compute-opt} (right), {\textsc BLT} models consistently perform on par with, or even surpass, their bpe counterparts, a pattern that persists across varying model scales and computational budgets. To our knowledge, {\textsc BLT} represents the first byte-level Transformer architecture to demonstrate comparable scaling behaviors to BPE-based models under compute-optimal settings. This observation supports our hypothesis that the ideal parameter-to-training compute ratio established for bpe models is likewise applicable to {\textsc BLT}, or at the very least, not significantly different.

Both enhancements to model architecture and the application of dynamic patching are essential for achieving the desired scaling behavior with respect to bpe. In ~\autoref{fig:scaling-compute-opt} (left), we present a comparison between space-patching-derived models and Llama 3. Here, SpaceByte \citep{slagle2024spacebyte} is estimated using {\textsc BLT} space-patching, deliberately omitting n-gram embeddings and cross-attention mechanisms. While SpaceByte provides better results compared to Megabyte, a significant gap remains when compared to Llama 3. The right side of \autoref{fig:scaling-compute-opt} demonstrates that advancements in both architectural design and dynamic patching yield notable performance gains. At scale, {\textsc BLT} approaches achieve results comparable to those of leading tokenizer-based systems, including Llama 3.

Furthermore, we find that tokenizer selection has a noticeable impact on the performance of tokenizer-based approaches. Specifically, models trained with the Llama-3 tokenizer exhibit superior results to those utilizing the Llama-2 tokenizer on identical datasets.

Ultimately, the {\textsc BLT} architecture demonstrates intermediate scaling behavior between Llama 2 and Llama 3 when operating with notably larger patch sizes. While the bpe tokenizers used in Llama 2 and Llama 3 yield average token lengths of 3.7 and 4.4 bytes, respectively, {\textsc BLT} is able to match these scaling characteristics at mean patch sizes of 6 or even 8 bytes. Given that inference flop is inversely related to the average patch size, utilizing an 8-byte patch size in {\textsc BLT} results in almost 50\% reduction in inference compute costs. Moreover, increasing the patch size appears to correlate with improved performance as both model and dataset sizes are scaled up. Notably, a {\textsc BLT} variant with an 8-byte patch size begins at a clear disadvantage compared to bpe Llama 2 at 1B parameters, but outperforms its bpe counterpart at 7B scale. This observation implies that such larger patch sizes could yield even greater benefits at higher model scales, and that utilizing even larger patches might become practicable as model and computational resources are further scaled.

\subsection{Beyond Compute Optimal Task Evaluations}
\label{section:evals}
To further investigate scaling dynamics, we train an 8B {\textsc BLT} model exceeding the compute-optimal threshold using the {\textsc BLT}-1T dataset, which provides both greater volume and improved quality. Model performance is then assessed on a diverse selection of classification and generation benchmarks. For these evaluations, we focus on tasks covering common sense reasoning, general world knowledge, and code generation abilities:
\paragraph{Classification tasks} encompass ARC-Easy (0-shot)~\citep{Eval_ARC}, Arc-Challenge (0-shot)~\citep{Eval_ARC}, HellaSwag (0-shot)~\citep{Eval_hellaswag}, PIQA (0-shot)~\citep{Eval_piqa}, and MMLU (5-shot)~\citep{hendrycks2020measuring}. Our approach employs prompt scoring, in which we determine the model's probability assigned to candidate answer choices, and we present mean accuracy results across these benchmarks.
\paragraph{Coding related generation tasks:}  To evaluate Python code synthesis, we utilize pass@1 on MBPP (3-shot)~\citep{Eval_MBPP} and HumanEval (0-shot)~\citep{Eval_HumanEval}, providing a direct assessment of model proficiency in coding scenarios.

In \autoref{tab:evals}, we present a comparative analysis of three models, each trained on the {\textsc BLT}-1T dataset: a bpe Llama 3 tokenizer-based model,\footnote{The Llama 3 tokenizer, featuring a 128k vocabulary, is selected over Llama 2’s 32k version due to its superior performance.} along with two variations of the {\textsc BLT} model. One variant, {\textsc BLT}-Space, implements a space-patching method, whereas the other, {\textsc BLT}-Entropy, adopts an entropy-driven patching strategy. This entropy model incorporates an approximate monotonicity constraint and resets its context upon encountering new lines, in accordance with the approach detailed in~\autoref{section:entropy-context}. Each model is trained utilizing a comparable amount of computational resources (flops). Notably, for {\textsc BLT}-Entropy, we decrease the entropy threshold used during inference from 0.6 to 0.1. This threshold reduction yields better performance across tasks, albeit with an increase in the number of inference steps required.

The {\textsc BLT}-Entropy model achieves superior results compared to the Llama 3 model on 4 of the 7 evaluated tasks, despite both models being trained on an equivalent data volume. These gains are likely attributable to two key factors: (1) more effective allocation of training resources through dynamic patching and (2) explicit modeling of information at the byte level, rather than at the token level.

Conversely, {\textsc BLT}-Space exhibits lower performance than the Llama 3 tokenizer across all tasks except one. However, it substantially decreases inference FLOPs owing to its greater mean patch size of 6 bytes. In contrast, both the bpe and entropy-patching models yield an average patch size close to 4.5 bytes on the training dataset mix. Given an equivalent training budget, the model utilizing the larger patch size is able to process 30\% more data relative to the other two models. This increased data coverage could potentially cause {\textsc BLT} to deviate further from the compute-optimal regime.

\subsection{Patches Scale Better Than Tokens} The {\textsc BLT} architecture permits joint scaling of both model capacity and patch size without exceeding a fixed compute budget for training and inference, nor requiring additional training data. Unlike fixed-vocabulary token-based approaches, patch-based models allow patch size to be increased without constraint—an advantage highlighted in Section~\ref{section:bpe-patch}. Utilizing larger patches reduces computational costs, enabling these saved resources to be invested in enlarging the global latent transformer, which benefits from being invoked less frequently.

A fixed inference scaling analysis is carried out to evaluate whether increasing model size and utilizing fewer steps on larger patches yields superior performance compared to deploying smaller models with more steps. We begin with models containing 400 million and 3.6 billion parameters using the Llama 2 tokenizer. Subsequently, we identify flop-equivalent models utilizing the Llama 3 tokenizer as well as {\textsc BLT}-Entropy models with mean patch sizes of 6 and 8 bytes, trained on the same data mix (refer to~\autoref{tab:fixed-inf-params} for the specifics of each model). For those models configured with an 8-byte patch size, three encoder layers are employed rather than a single layer. Each model undergoes training under a range of specified training flop budgets.

\autoref{fig:fixed_inference_scaling} illustrates that {\textsc BLT} models exhibit superior scaling behavior compared to tokenization-based architectures under both inference flop constraints. While BPE models tend to excel when training resources are limited, {\textsc BLT} quickly overtakes them once moving beyond the compute-optimal threshold. In practical terms, investing additional resources during the initial pretraining phase can lead to a higher performing model even when inference costs are held constant. A notable case in point is the 8B model class, exemplified by Llama 3.1, which was pretrained on data volumes two orders of magnitude larger than the compute-optimal amount for that model size.

The point at which {\textsc BLT} begins to surpass token-based models has moved somewhat nearer to the compute-optimal threshold with the adoption of models from higher flop classes, shifting from 3x to 2.5x the compute-optimal allocation. Likewise, the model utilizing a larger patch size of 8 demonstrates a more pronounced improvement in scaling for the higher flop category, surpassing alternative models earlier. As elaborated in Section~\ref{section:parameter-matched-scaling}, at larger model sizes, using larger patch sizes tends to yield performance more comparable to BPE models. We partly attribute this effect to the diminishing proportion of total computation devoted to the byte-level Encoder and Decoder modules, which appear to exhibit slower scaling relative to the Latent Transformer. Notably, when total parameter count increases 20-fold, from 400M up to 8B, the local model parameters for {\textsc BLT} only increase by about 2x. This observation is crucial, since only the Latent Transformer with larger patch sizes impacts computational requirements, leaving the byte-level modules unchanged. Indeed, this explains why {\textsc BLT}-Entropy with patch size 8 increased from 1.6x to 1.7x the parameter count of Llama 2 when transitioning to the higher model scale.

To summarize, our investigation into patch-length scaling reveals that the {\textsc BLT} patch-based framework attains superior scaling behavior when both the size of patches and the overall model are enlarged together. These advantageous scaling patterns appear to hold and potentially strengthen with increasing model size.

\section{Byte Modeling Improves Robustness} We further assess the robustness of {\textsc BLT} in comparison to token-based counterparts that do not directly incorporate byte-level data, and introduce a method for converting pretrained token models to operate at the byte level.
\label{sec:robustness}

\subsection{Character-Level Tasks}

An initial impetus for developing byte-level models stemmed from their resilience to noise present at the byte level within inputs and their capability to recognize the substructure of tokens—a challenge that tokenizer-based models often face. To assess these aspects, we conduct supplementary experiments using benchmark datasets specifically designed to gauge models’ robustness against input noise as well as their sensitivity to character-level information, spanning both English and multilingual text, and covering elements such as digits and phonemes. The findings from these evaluations are provided in Table~\ref{tab:char_tasks}.

\paragraph{Noisy Data} To assess model robustness, we generate altered versions of the benchmark classification tasks outlined in Section~\ref{section:evals}, allowing for direct comparison between tokenizer-based models and {\textsc BLT}. Specifically, we apply five distinct character-level perturbation techniques to modify the input data: (a) \textit{AntSpeak}: Reformats the entire text into space-separated uppercase letters; (b) \textit{Drop}: Randomly deletes 10\% of characters within the text; (c) \textit{RandomCase}: Randomly assigns uppercase or lowercase to each character, resulting in 50\% of characters being in each case; (d) \textit{Repeat}: Selects 20\% of characters and duplicates them up to four times; (e) \textit{UpperCase}: Converts every character in the text to uppercase. For evaluation, each perturbation is applied separately to the prompt, the completion, or both, and the mean performance across these tasks is calculated. Table~\ref{tab:char_tasks} displays the outcomes for the noised HellaSwag~\citep{Eval_hellaswag} data, where we observe that {\textsc BLT} consistently demonstrates superior robustness compared to tokenizer-based architectures, exhibiting an average performance margin of 8 points against a similarly-trained model, and even outperforming the Llama 3.1 variant that leverages a substantially larger training corpus.

\paragraph{Phonology - Grapheme-to-Phoneme (G2P)} We evaluate how effectively {\textsc BLT} converts sequences of graphemes (the letters constituting a word) into their corresponding phonemic transcriptions (representing pronunciation). As shown in Table \ref{tab:char_tasks}, results from the G2P task—performed under a 5-shot learning paradigm with Phonology Bench~\citep{suvarna2024phonologybench}—demonstrate that {\textsc BLT} surpasses the Llama 3 1T tokenizer-based baseline model on this metric.

\paragraph{CUTE}
In order to evaluate character-level comprehension, we test {\textsc BLT} using the CUTE benchmark~\citep{edman2024cute}. This benchmark encompasses a range of tasks divided into three main areas: compositional understanding, recognition of orthographic similarity, and sequence manipulation skills. The CUTE suite poses a notable difficulty for typical tokenizer-based architectures, which tend to be aware of token spellings, yet often fall short in leveraging this information effectively for text manipulation. As reported in Table~\ref{tab:char_tasks}, {\textsc BLT}-Entropy surpasses both BPE Llama 3 variants by more than 25 points across the benchmark. Notably, the model displays outstanding performance on tasks requiring character manipulation, achieving 99.9\% accuracy on both spelling challenges. These substantial gains, especially given that {\textsc BLT} is trained on sixteen times less data than Llama 3.1, suggest that learning character-level features remains a significant hurdle for BPE-based approaches. Figure \ref{fig:cute_examples} provides illustrative examples highlighting cases where Llama 3 tokenizer models encounter difficulties, while our {\textsc BLT} system excels. The only tasks where BPE shows superior results are word deletion and insertion. Although manipulating words at the byte-level may pose challenges, the performance gap is relatively small, and it may ultimately be more tractable to compose words from characters than to infer character-level details from word-level representations. Our experiments utilize consistent evaluation procedures and employ the original Huggingface prompts for all tasks. It is possible that BPE models could see improved results with the application of additional prompt engineering.

\paragraph{Low Resource Machine Translation} We assess {\textsc BLT} on translation tasks involving twenty-one low-resource languages spanning six widely recognized language families, featuring diverse scripts, using the FLORES-101 benchmark~\citep{goyal2022flores}. The corresponding SentencePiece BLEU scores are presented in Table~\ref{tab:flores}. Our findings show that {\textsc BLT} delivers a 2-point overall improvement when translating into English and a 0.5-point gain when translating from English, relative to a model based on the Llama 3 tokenizer. For high-frequency language pairs, the performance of {\textsc BLT} is at least on par with, or marginally better than, that of Llama 3. Notably, across a wide range of lower-resource language family pairs, {\textsc BLT} consistently surpasses Llama 3, highlighting byte modeling's capability to generalize more effectively to rare and diverse byte sequence patterns.

\subsection{Training {\textsc BLT} from Llama 3}
\label{sec:distillation}
We investigate a method whereby {\textsc BLT} models can benefit from pre-trained tokenizer-based models, with the aim of accelerating and enhancing training convergence. Specifically, the global transformer parameters in {\textsc BLT} are initialized with those from a pre-trained Llama 3.1 model. During subsequent training, the global transformer weights are updated using a learning rate that is one-tenth of what is used for the local encoder and decoder. Training is conducted for the optimal number of steps as recommended for Llama 3, and results are compared against a baseline {\textsc BLT} in Table~\ref{tab:distill}. The findings indicate that initializing {\textsc BLT} from Llama 3.1 yields performance that notably exceeds both the Llama 3 and the baseline {\textsc BLT}, with all models being trained under an equivalent computational budget. Furthermore, even when compared against the {\textsc BLT}-Entropy model (refer to Table~\ref{tab:evals})—which was exposed to a much larger dataset of 1T tokens—{\textsc BLT} initialized from Llama 3.1 demonstrates better performance on the MMLU benchmark. This suggests that such initialization may offer a highly efficient path for reducing the required training flops.

This configuration enables the conversion of tokenizer-based architectures into tokenizer-free models, thereby allowing a pre-trained LLaMA 3.1 model to serve as a {\textsc BLT} model. For a thorough evaluation, we report results for the original LLaMA 3.1, trained on 15T tokens, alongside those for the LLaMA-derived {\textsc BLT} in Table \ref{tab:distill}. Although our approach demonstrates only minor degradation in performance on MMLU and HumanEval, it incurs notably greater losses on additional benchmarks. These findings indicate that additional research is necessary to maximize the benefits of pre-trained models—specifically, to enhance performance via more effective tuning of data mixtures and other relevant hyperparameters.

\section{Ablations and Discussion}
\label{section:ablations}
This section presents ablation studies aimed at validating design decisions in the architecture of {\textsc BLT}, as well as examining the patching strategy and hyper-parameter selections for the 8B parameter {\textsc BLT} model trained on the {\textsc BLT}-1T dataset.

\paragraph{Entropy Model Hyper-parameters} To evaluate the influence of entropy model size and the length of the context window on model scaling behavior, we train byte-level entropy transformer models spanning parameter counts from 1m to 100m, employing context window sizes that vary between 64 and 512. The scaling laws relating bpb to training FLOPs are shown using our $400m$ and $1b$ {\textsc BLT} models trained on Llama-2 data and are illustrated in \autoref{fig:entropy_ablation}. Our observations indicate that both increasing the entropy model's size and its context window length tend to enhance scaling performance, but the improvements become marginal for model sizes exceeding 50m parameters.

\paragraph{Types of Patching}

We conduct an ablation study on the four patching methods described in Section~\ref{section:patching}: 1) Strided Patching, tested with strides of 4 and 6, 2) Whitespace-based Patching, 3) Patching with BPE Tokenizer according to the Llama 3 tokenizer, and 4) Entropy-driven patching utilizing a compact byte-level language model.

Although dynamic patching shortens the actual sequence length, we standardize sequence lengths across all patching methods to ensure comparable context windows. This approach guarantees that, on average, each model is exposed to the same number of bytes per sequence during both training and inference, thus eliminating the potential for confounding due to varying context sizes. The outcomes of these controlled ablations are presented in Figure~\ref{fig:scaling-compute-opt}. All alternative patching strategies demonstrate superior performance relative to static patching, and space patching, in particular, performs nearly as well as dynamic entropy-based patching.

In \autoref{tab:evals_ablation}, we report the results of benchmarking {\textsc BLT} models that utilize either tokenizer-based approaches, space patching, or entropy-based patching. All models were trained on the Llama 2 dataset, using an optimal number of training steps as outlined in~\citep{dubey2024llama}. While space patching represents a more straightforward method, as it does not require executing an entropy model dynamically during training, our results indicate that the improvements from entropy-based patching identified in scaling experiments~(Section~\ref{section:scaling}) are also observed when assessing performance on downstream benchmark tasks.\footnote{The space patching results were produced from earlier training runs without cross-attention; however, similar outcomes are evident with cross-attention enabled.}

\paragraph{Cross-Attention} 
In~\autoref{tab:crossattn_ablations_1b}, we systematically assess the impact of introducing cross-attention at different locations within the encoder and decoder architectures of {\textsc BLT}. For cross-attention in the encoder, we compare several query initializations: (1) using a uniform learned embedding across all global states, (2) employing a hash embedding derived from the bytes contained in the patch, and (3) applying pooling over the encoder’s hidden states corresponding to the patch bytes at the specified encoder layer.

Our results indicate that incorporating cross-attention within the \textit{decoder} yields the greatest benefits. Applying cross-attention in the encoder produces only minor gains, and these occur exclusively when query initialization is performed via pooling. Furthermore, our analysis reveals that cross-attention offers notable advantages on Common-Crawl datasets, with the improvement being even more pronounced as patch sizes increase.

\paragraph{n-gram Hash Embeddings} 

We investigate the effect of varying n-gram hash embedding vocabulary sizes—specifically 0, 100K, 200K, and 400K—and report the findings in Table~\ref{tab:ngram_ablations_1b}. Our analysis reveals that hash embeddings yield improvements across all examined domains, with particularly notable benefits for Wikipedia and Github datasets (a 0.04 bpb improvement, compared to a 0.01 bpb gain after 15k steps at the 8B model scale). At the 8B scale, reducing the number of hashes from 500K to 300K results in only a marginal performance drop of approximately 0.001 bpb at 15k steps. These results suggest that hash embeddings play a crucial role in aligning the performance of {\textsc BLT} with that of models relying on tokenizers, but the marginal improvements decrease once the hash count surpasses 300K. Furthermore, the observed benefits of hash embeddings appear to supplement those offered by cross-attention, as both features yield distinct improvements on various datasets.

\paragraph{Local Model Hyperparameters} Table \ref{tab:local_layers} presents an ablation study of different configurations regarding the number of layers utilized in the local encoder and decoder. The results indicate that {\textsc BLT}, when used alongside hash n-gram embeddings, performs effectively with a highly minimal encoder—specifically, one consisting of a single layer—while relying on a comparatively deeper decoder.

\section{Related Work}
\label{section:related_work}

\textbf{Character-Level RNNs:} Since the advent of neural architectures, Character Language Modeling has gained significant traction as a research area~\citep{sutskever2011generating,mikolov2012subword,graves2013generating}, primarily due to its inherent capability to naturally handle out-of-vocabulary terms, bypassing the need for traditional back-off techniques. In their work, \cite{kim2016character} construct a model that exclusively ingests character sequences on the input side via convolutional and highway layers, which subsequently interface with LSTM-based RNNs. This approach achieved results that rivaled contemporaneous state-of-the-art RNN language models for English, and even surpassed them in languages with complex morphology—a notable strength of character-level LLMs. Building on this, \cite{kenter2018byte} demonstrate that byte-level LSTM architectures excel over word-level counterparts in machine comprehension tasks, particularly for morphologically complex languages such as Turkish and Russian. Along similar trajectories, \cite{zhang2015character} implement character-based convolutional neural networks for text classification, achieving superior performance compared to word-based models on select tasks. Further advances include the hierarchical LSTM model proposed by \citet{chung2019hierarchical}, which employs boundary detectors to infer latent structural hierarchies in text and thereby bolsters character-level language modeling results. Departing from attention-based schemes, ByteNet by \cite{kalchbrenner2016neural} employs convolutional layers directly on character sequences for machine translation tasks.

\textbf{Character-Level Transformers:} The introduction of attention-based transformer architectures~\citep{vaswani2017attention} alongside subword tokenization strategies~\citep{sennrich-etal-2016-neural} markedly advanced the capabilities of neural models in language modeling tasks and standardized benchmarks. Despite these gains, utilizing word and sub-word tokens inherently establishes an inductive bias about the preferred abstraction layer for model processing. Seeking to bridge transformer effectiveness with early success in character-level language modeling, \citet{al2019character} implemented highly deep transformer networks. By leveraging auxiliary loss objectives, they managed to train character-level transformer models that surpassed the performance of earlier character LSTM models. Nonetheless, a notable disparity persisted compared to word-level large language models. Furthermore, GPT-2~\citep{radfordgpt2} documented that, on expansive datasets such as the 1 billion word benchmark, language models operating at the byte level did not achieve competitive results relative to their word-level counterparts.

According to \cite{Choe2019BridgingTG}, transformer-based LLMs that operate at the byte level can achieve superior performance compared to their subword-based counterparts with similar model sizes. However, these byte-level models require significantly more computational resources and longer training durations. In a related vein, \cite{el2020characterbert} present CharFormer, a variant of BERT that generates word representations via convolutions over character embeddings, and report performance gains within the medical domain, albeit with substantially increased computational demands. The CANINE model introduced by \cite{clark2022canine} features 150M parameters and processes raw character sequences. At its core is a deep transformer, resembling the architecture of our global model, accompanied by a local transformer module and strided convolution layers to reduce input character resolution. CANINE achieves superior results compared to the similarly sized token-level mBERT on various multilingual benchmarks. ByT5~\citep{xue2022byt5} proposes byte-level encoder-decoder architectures that intentionally avoid patching mechanisms. While their approach enhances noise robustness and achieves competitive performance relative to tokenizer-based systems using only one-fourth the amount of training data, the absence of patching requires costly attention computations for each input byte, resulting in substantial computational overhead. Representing sequences as bytes, instead of subword tokens, further amplifies input sequence length, making the scaling of byte-level models computationally challenging. More recently, utilizing the Mamba Architecture~\citep{gu2023mamba}, which supports maintaining a fixed-size state across extensive context windows, \cite{wang2024mambabyte} train byte-level Mamba models—again forgoing patching—and demonstrate that they outperform comparable byte-level transformer architectures under FLOP-matched constraints at the 350M parameter scale, as measured by bits-per-byte across multiple datasets.

\textbf{Patching-based approaches:} Patching strategies offer a promising way to reduce the substantial computational overhead commonly associated with byte-level LLMs, all while potentially preserving performance levels. Several studies have reported early-stage success with such methods when applied to models of modest size and datasets with relatively low numbers of training bytes. \cite{nawrot-etal-2022-hierarchical} introduced static patching via downsampling and upsampling and proposed the hourglass transformer architecture, which achieves superior results compared to existing byte-level approaches at the 150M parameter scale. Building upon this work, \cite{nawrot-etal-2023-efficient} incorporate dynamic patching mechanisms, including an end-to-end trainable boundary-predictor, a supervised boundary-predictor leveraging specific tokenizers, and an entropy-guided patching approach akin to {\textsc BLT}. Their results demonstrate performance gains over contemporaneous vanilla transformers in language modeling at the 40M parameter range, evaluated on roughly 400M tokens. In related work, \cite{lester2024training} explore the use of sequences compressed via arithmetic coding to attain compression ratios surpassing those offered by BPE. By introducing an equal-info windows technique, they achieve improved outcomes over byte-level baselines for language modeling, although performance still lags behind subword-based models.

Our study is primarily influenced by, and most comparable to, MegaByte~\citep{yu2023megabyte}. MegaByte is a decoder-only, causal large language model that applies a predetermined, unchanging strategy to form patches from bytes by concatenating their representations, subsequently employing a local decoder-side model to map patches back to byte sequences. Their findings indicate that, at the 1B parameter scale and on datasets totaling roughly 400B bytes, MegaByte attains performance on par with models built around tokenizers. In our experiments, we include MegaByte in all comparative analyses, observing that static patching falls short when compared to present state-of-the-art, compute-optimized tokenizer-based architectures under fixed computational budgets. We further show that {\textsc BLT} narrows this performance gap. Similarly, \cite{slagle2024spacebyte} highlight the same limitations in MegaByte, proposing modifications such as patching on whitespace and related byte categories, and introducing a local encoder component. Their results suggest certain advantages over tokenized transformer models on domains like Github and arXiv at the 1B parameter scale within compute-constrained scenarios. We extend the experimental evaluation to include these models, and our evidence suggests that additional architectural innovations are necessary for byte-level approaches to achieve scalability and to rival top-performing token-based models, including Llama 3.

\section{Limitations and Future Work}

For the purposes of making architectural decisions in this study, we train models using the optimal number of steps as established for Llama~3~\citep{dubey2024llama}. Nevertheless, these scaling laws were originally derived for transformers utilizing BPE, and may not yield ideal (data, parameter) ratios for {\textsc BLT}. Determining specific scaling laws tailored for {\textsc BLT} is left for future research, as these could reveal more advantageous scaling patterns for our model. Furthermore, most of the experiments herein are conducted with models up to 1B parameters. There remains a possibility that different architectural choices will prove more effective as model sizes grow to 8B parameters and larger, potentially resulting in enhanced performance at greater scales.

Current transformer codebases and libraries are primarily optimized for architectures utilizing tokenizers. Although we conduct experiments matching theoretical flops and employ efficient solutions—like FlexAttention—for non-standard transformer layers, our implementations might still lag behind tokenizer-based models in actual runtime performance and could require additional optimizations to reach comparable speed.

Whereas {\textsc BLT} relies on a separately pre-trained entropy model to facilitate patching, integrating the learning of the patching mechanism into an end-to-end framework presents a promising avenue for future investigation. In Section \ref{sec:distillation}, we provide preliminary experimental evidence suggesting that it is feasible to "byte-ify" tokenizer-based models—such as Llama 3, which are trained on datasets exceeding 10T tokens—by transferring their weights to initialize the global transformer component and keeping these weights fixed. Further exploration in this area could yield approaches that preserve the advantages of bytefying while potentially enhancing model performance beyond that of the original tokenizer-dependent architectures, all without necessitating complete retraining.

\section{Conclusion}
\label{section:conclusion}

In this work, we introduce the Byte Latent Transformer (\textbf{BLT}), a novel model architecture that challenges the standard reliance on fixed-vocabulary tokenization within large language models. {\textsc BLT} incorporates a flexible, learnable byte-grouping mechanism, assembling bytes into patches in a way that dynamically adjusts computational effort according to input complexity. This enables substantial gains in model efficiency and resilience. Through a comprehensive scaling analysis, we show that {\textsc BLT} attains performance parity with tokenization-based architectures such as Llama 3 at sizes up to 8B parameters and 4T bytes, and offers the opportunity to achieve up to 50\% fewer inference floating point operations with only marginal declines in evaluation results. In addition, {\textsc BLT} allows the scaling of both the underlying model and the patch size concurrently within a predetermined inference computation budget—introducing an advantageous scaling axis for real-world compute scenarios. By processing data directly in its raw byte form, {\textsc BLT} enhances robustness to data irregularities and offers improved modeling of rare or complex sub-word features. Collectively, our findings establish {\textsc BLT} as a compelling substitute for fixed-token approaches, delivering an efficient, scalable, and robust architecture suited for next-generation adaptive language models.

\end{document}